% Created 2026-03-07 Sat 16:58
% Intended LaTeX compiler: xelatex
\documentclass[11pt]{article}
\usepackage{graphicx}
\usepackage{longtable}
\usepackage{wrapfig}
\usepackage{rotating}
\usepackage[normalem]{ulem}
\usepackage{capt-of}
\usepackage{hyperref}
\usepackage{geometry}
\geometry{margin=1in}
\usepackage{setspace}
\usepackage{amsmath}
\usepackage{amssymb}
\usepackage{booktabs}
\onehalfspacing
\author{Bradley K. DePuy, Independent Researcher}
\date{2026}
\title{The Medium — A Foundational Theory of Reality}
\hypersetup{
 pdfauthor={Bradley K. DePuy, Independent Researcher},
 pdftitle={The Medium — A Foundational Theory of Reality},
 pdfkeywords={},
 pdfsubject={},
 pdfcreator={Emacs 29.3 (Org mode 9.7.39)}, 
 pdflang={English}}
\begin{document}

\maketitle
\tableofcontents

Version 0.2 | DOI: 10.5281/zenodo.18804587 | CC BY 4.0

\emph{Out of empty chaos all of creation becomes.}

That sentence is the theory. Everything that follows is the attempt to say it precisely.
\section{The Foundational Claim}
\label{sec:orgd40b3c3}

The framework rests on a single foundational claim from which everything else follows:

\begin{quote}
The Medium always is and always was. It has no beginning and no end. It exists
outside of any temporal framework. Intrinsic distinction is an eternal property
of the Medium — not an event that occurred but a condition that always obtains.
The structure we observe is not the result of something that happened. It is the
expression of what the Medium always is.
\end{quote}

This claim does three things at once.

It eliminates the need for a trigger. The question of what initiated the universe
only arises if time is assumed to be primitive. If the Medium always is and always
was, there is no before and no initiation. The trigger problem dissolves not by
answering it but by identifying it as a consequence of a false assumption.

It eliminates time as a primitive. The Medium exists outside of time. What physics
calls time is not a property of the substrate. It is an experience of observers
embedded within the differentiated structure of the Medium — a consequence of
transition, not a fundamental feature of reality.

It reframes differentiation. The emergence of structure from the Medium is not a
process that started. It is a condition that always obtains. The ground state and
the differentiated state are not before and after. They coexist within the Medium
eternally, because intrinsic distinction — \(\iota\) — is eternal, and time does
not apply to it.
\section{The Stochastic Character of the Medium}
\label{sec:org2d542c9}

Before the formal axioms are introduced, the nature of the Medium's ground state
requires precise characterization. That characterization is stochastic — and the
stochastic framing is not incidental. It is the most exact description available
of what the Medium is prior to the activation of intrinsic distinction.
\subsection{The Ground State as Fully Random Autonomous Stochastic System}
\label{sec:orgbe8bf32}

The Medium at its ground state is a fully autonomous stochastic system. It
requires no inputs, no external cause, no temporal framework, and no governing
distribution. Its randomness is not randomness within a framework — it is
randomness without framework. No outcome is privileged over any other. No
state is more probable than any other. Non-existence holds no advantage over
existence. Existence holds no advantage over non-existence.

This is not chaos. Chaos retains structure. The ground state of the Medium
precedes structure entirely. It is the absolute suspension of all constraint,
including constraints upon randomness itself. It is meta-randomness — randomness
that is random about its own nature.
\subsection{What Fully Random Stochastic Nothingness Implies}
\label{sec:org553a6a1}

In a fully random autonomous system of infinite extent with no preferred states,
every possible configuration has nonzero probability. This has a consequence that
is not intuitive but is mathematically unavoidable:

\begin{quote}
In an infinite substrate with fully random stochastic character and no privileged
states, every possible fluctuation — including the emergence of ordered, structured,
low-entropy configurations — is not merely possible but necessary as a feature of
the distribution itself.
\end{quote}

The universe is not something that happened against the odds. It is a necessary
expression of a substrate that has no mechanism for suppressing any possibility.
The question \emph{why something rather than nothing} dissolves: it assumes nothingness
is the stable default. In a fully random stochastic substrate, no state is default.
No state is stable. No state is privileged.
\subsection{The Relationship Between Stochastic Character and \(\iota\)}
\label{sec:org028187b}

The stochastic character of the Medium describes the nature of the ground state.
Intrinsic distinction — \(\iota\) — describes the property that makes the ground
state generative rather than merely static.

Without \(\iota\), the stochastic ground state remains flat and undifferentiated.
All possibilities are present but none are actualized. Pure potential with no
capacity for manifestation.

With \(\iota\), the fully random substrate has the capacity to produce local,
spontaneous reductions in entropy — local distinctions, local structure, local
reality. \(\iota\) is what makes the randomness productive. It is the one property
that transforms infinite stochastic potential into actual differentiated existence.

The relationship is precise:

\begin{itemize}
\item Stochastic character describes what the ground state is
\item \(\iota\) describes what makes the ground state capable of producing structure
\item Manifestation \(V\) is the stochastic event — the actualization of a possibility
\item Structure formation at \(c\) is what a successful manifestation produces
\end{itemize}
\section{Axioms and Primitives}
\label{sec:orga316e64}

The formal definitions follow directly from the foundational claim and the
stochastic characterization of the ground state. Several are marked provisional.
Open questions are named precisely and carried explicitly as objects of
investigation rather than concealed as weaknesses.
\subsection{Axiom I — The Eternal Medium}
\label{sec:org9383354}

The Medium always is and always was. It has no beginning, no end, and no external
cause. It is not an event. It is the substrate within which all differentiated
structure occurs. There is no moment at which the Medium came into existence and
no moment at which it ceases.

\begin{equation}
\nexists\, t_0 : \mathcal{M} \text{ begins at } t_0 \qquad
\nexists\, t_1 : \mathcal{M} \text{ ends at } t_1
\end{equation}

\emph{Status: Axiomatic.}
\subsection{Axiom II — Intrinsic Distinction}
\label{sec:org77a1eb1}

The Medium possesses intrinsic distinction. This is not a process that occurs,
not an event that happened, and not a property imposed from without. It is what
the Medium is.

This property is denoted \(\iota\). It is not an operator, not a function, not a
process. It is the one irreducible property of the thing that simply is — the
capacity of the undifferentiated to give rise to the differentiated. Its precise
nature is the central investigation of the theory.

The activation of \(\iota\) — the event of intrinsic distinction becoming manifest
— is denoted \(V\). \(V\) is what occurs when the stochastic potential of the Medium
is actualized at a location. It is the stochastic event described in Section II.

\emph{Status: Axiomatic. Nature of \(\iota\) under investigation.}
\subsection{Axiom III — Locality Emergence}
\label{sec:org51f8f99}

Location is not a primitive of the Medium. It does not exist prior to
differentiation. Location is produced when differentiation, propagating at the
limit \(c\), gives rise to self-sustaining structure. Before this condition is met,
the concepts of here and there, near and far, position and distance are genuinely
absent — not suppressed but nonexistent.

\emph{Status: Axiomatic.}
\subsection{Primitive I — The Medium}
\label{sec:org0e76341}

In its first formal approximation, valid only after locality has been established
per Axiom III:

\begin{equation}
M = \{x \mid x \text{ is a location}\}
\end{equation}

The Medium prior to locality resists set-theoretic description. It is
pre-extensional: it has no members because it has no parts. It simply is,
and \(\iota\) is what it is.

\emph{Status: Provisional — applies post-locality only.}
\subsection{Primitive II — The Population Function}
\label{sec:org78f4954}

The population function \(\phi\) records where and when \(V\) — the activation of
\(\iota\) — has occurred. It is a partial function defined only at locations where
manifestation has taken place.

\begin{equation}
\phi : M \rightarrow V \qquad \phi(x) \text{ exists iff } x \text{ has manifested}
\end{equation}

\emph{Status: Current.}
\subsection{Primitive III — The Manifestation Probability}
\label{sec:org9ebb1fa}

Manifestation is a binary stochastic event. At any transition \(t\), a location
\(x \in M\) either manifests or it does not. The stochastic character of the
ground state means that no location is guaranteed to manifest and none is
guaranteed not to:

\begin{equation}
\text{manifestation}(x,t) = \begin{cases} 1 & \text{if } x \text{ manifests at } t \\ 0 & \text{if } x \text{ does not manifest at } t \end{cases}
\end{equation}

\begin{equation}
p(x) = \Pr(\text{manifestation}(x,t) = 1) \qquad p : M \rightarrow [0,1]
\end{equation}

\emph{Status: Current.}
\subsection{Primitive IV — Transition}
\label{sec:org3035582}

Time is not a primitive of this framework. What physics describes as temporal
evolution is, at the level of the Medium, transition — a change from one state
of differentiation to another. The transition index \(t\) is a label, not a clock.
No duration is implied. No direction is implied. No causal relationship between
states is assumed.

\begin{equation}
M_t \rightarrow M_{t+1}
\end{equation}

\emph{Status: Current.}
\subsection{Primitive V — The Propagation Limit}
\label{sec:org6a3bb81}

The Medium imposes a fundamental constraint on how rapidly the activation of
\(\iota\) can propagate across a transition, denoted \(c\):

\begin{equation}
|\phi_{t+1}(x) - \phi_t(x)| \leq c
\end{equation}

\(c\) precedes light, matter, and all physical phenomena. The speed of light is
proposed to be a consequence of \(c\), not its definition.

\emph{Status: Provisional.}
\subsection{Primitive VI — The Structure Formation Condition}
\label{sec:org683c62b}

A self-sustaining structure forms when propagation reaches the limit exactly:

\begin{equation}
|\phi_{t+1}(x) - \phi_t(x)| = c
\end{equation}

This is the threshold at which locality is born per Axiom III. Space does not
precede this structure — space is what this structure makes possible. Everything
we call physical reality is proposed to be the behavior of such structures and
their interactions within the Medium.

\emph{Status: Provisional.}
\section{The Entropic Arc}
\label{sec:org86f55e4}

The relationship between the stochastic ground state of the Medium and the
thermodynamic concept of entropy yields a consequence that bears directly on
one of the deepest unsolved problems in the current understanding of physical
reality.
\subsection{The Ground State as Maximum Entropy}
\label{sec:org5463a77}

Entropy measures the number of possible arrangements of a system consistent with
its observed condition. The more arrangements that produce the same observable
state, the higher the entropy of that state.

The ground state of the Medium — fully random, fully undifferentiated, no
resolved structure, no preferred states — is a maximum entropy state by this
definition. There is no arrangement distinguishable from any other. Every
configuration of the substrate is equivalent. The number of arrangements
consistent with the ground state is maximal. Entropy is maximal.

The stochastic character of the ground state and its maximum entropy character
are the same condition described in two different formal languages:

\begin{equation}
S = k \log W \qquad \text{where } W \text{ is maximal at the ground state}
\end{equation}
\subsection{Intrinsic Distinction as a Low Entropy Event}
\label{sec:org6cab2f4}

The activation of \(\iota\) — the first stochastic manifestation of distinction
in the maximum entropy ground state — is necessarily a minimum local entropy
event. One specific structure emerges from an infinite sea of equally probable
arrangements. The number of arrangements consistent with this outcome is
vanishingly small relative to the total. This is minimum local entropy by
definition.
\subsection{Resolution of the Low Entropy Initial Condition Problem}
\label{sec:org08c07da}

The current understanding of physical reality faces a foundational problem that
has resisted resolution: the universe appears to have begun in a state of
extraordinarily low entropy. Nothing in the existing physical framework explains
this. It sits as a brute unexplained initial condition at the foundation of the
standard model. The improbability of this initial state has been characterized
as staggering beyond ordinary description.

The Medium framework resolves this without being asked.

The low entropy initial condition of the universe is the direct and necessary
consequence of \(\iota\). The first distinction is a low entropy event because it
is the emergence of one specific structure from a maximum entropy stochastic
substrate. No separate explanation is required. The initial condition is not a
brute fact — it follows necessarily from the single primitive property of the
Medium combined with its stochastic character.

The framework therefore transforms one of the deepest unexplained facts in the
current understanding of physical reality into a necessary consequence of
Axiom II and Section II.
\subsection{The Full Entropic Arc}
\label{sec:orgac59f12}

The relationship between the stochastic ground state, \(\iota\), and entropy
implies a complete arc for physical reality within the Medium:

\begin{enumerate}
\item The Medium obtains eternally at maximum entropy — fully random, fully
undifferentiated, stochastic potential without actualization.
\item \(\iota\) manifests — the stochastic event actualizes — first distinction —
minimum local entropy — maximum information emergence.
\item Structure forms at \(c\) — locality is born — complexity builds — entropy
begins its ascent through successive transitions.
\item Stars, galaxies, and life emerge as temporary low entropy configurations —
eddies of order in the entropic stream.
\item Consciousness appears — the Medium becoming locally aware of its own
stochastic origin and entropic arc.
\item The trajectory continues toward maximum entropy — maximum undifferentiation
— the stochastic ground state of the Medium.
\end{enumerate}

The universe does not move from nothing to something and back to nothing. It
moves from the stochastic ground state of the Medium, through a single low
entropy fluctuation initiated by \(\iota\), and back to the ground state. Because
the Medium is eternal and \(\iota\) is its eternal property, this arc does not
occur once. It occurs wherever and whenever \(\iota\) manifests — infinite
stochastic fluctuations, infinite arcs, each spending its low entropy initial
condition and returning to the ground.

Not a beginning and an end. A breathing.
\subsection{A Provisional Characterization of \(\iota\)}
\label{sec:orgb18b6e1}

The entropic and stochastic framing together suggest a provisional formal
characterization of intrinsic distinction that avoids defining it in terms of
the differentiated system it produces:

\begin{quote}
Intrinsic distinction \(\iota\) is the property of the fully random maximum
entropy substrate that permits local, spontaneous entropy reduction — the
appearance of distinguishable structure in a stochastic ground state where
all arrangements are equally probable and no outcome is suppressed.
\end{quote}

This characterization connects \(\iota\) to thermodynamics and the stochastic
framing without circularity. Whether it is formally adequate or requires
genuinely new mathematical formalism is the central open question of the theory.
\section{The Problem With Where We Started}
\label{sec:org6f8a951}

Modern physics is extraordinary. Its achievements are beyond dispute. But the
foundations were built from the wrong starting point.

Human beings began making sense of the world from the human scale. That is
natural and inevitable. But it means the conceptual framework underlying physics
was constructed from the middle outward, not from the beginning. We built upward
toward the very large and downward toward the very small from a foundation that
was never examined as a foundation. It was simply where we happened to be
standing.

The consequences are visible at the edges of current theory.

The two most precisely verified theories in the history of science are
fundamentally incompatible with each other. At the scale of the very small,
physics works one way. At the scale of the very large, it works a different way.
Both cannot be fully correct. Both are approximations of something more
fundamental that has not yet been found. This is not a peripheral problem — it
is a contradiction between the two load-bearing pillars of modern physics.

The standard model of cosmology requires the entire universe to have emerged from
a singularity — a point at which the known laws of physics break down entirely.
It then requires unexplained phenomena with no known physical basis to account
for observations that do not fit the prediction. Approximately 95\% of the
universe by current models consists of placeholder labels for things that do not
fit. The practitioners of the field know this and say so.

The fundamental equations of physics are time-symmetric — they work equally well
forward and backward. The experienced directionality of time has no explanation
within the formal system. The extraordinarily low entropy initial condition of the
universe has no explanation within the formal system.

These are not minor technical problems awaiting refinement. They are foundational
cracks. They are consistent with what one would expect to find if time and
locality are not primitive but emergent — and if the substrate is stochastic
rather than deterministic at its deepest level.

The root of these problems is the same: time is assumed to be primitive. It is
not. Locality is assumed to be primitive. It is not. The ground state is assumed
to be lawfully constrained. It is not. Removing those assumptions does not
invalidate the achievements of modern physics. It reveals the substrate from
which those achievements can be derived — and from which the gaps that the field
itself has named may find resolution.
\section{What the Framework Proposes}
\label{sec:org7a7889e}

\subsection{The Medium}
\label{sec:org149837f}

The Medium is the primitive substrate of reality. At its most primitive level it
is one undifferentiated thing — without parts, without position, without internal
distinction of any kind. Its character is fully random and autonomous — the
stochastic ground state described in Section II. It is not a space containing
things. It simply is, and \(\iota\) — intrinsic distinction — is what it is.
\subsection{Nothing Moves}
\label{sec:org5d01081}

Nothing in the Medium moves. Movement requires something to move from one
location to another. But location is not primitive — it is an emergent consequence
of structure formation at \(c\). What we experience as motion is pattern — the
pattern of differentiation shifting across the Medium. Consider an LED signboard
displaying a moving message. The lights do not move. What moves is the pattern
of illumination. The message is real. The movement is not. Nothing traveled.
What we observe as a photon traveling from a distant star is the pattern of
differentiation shifting across the one undifferentiated stochastic substrate.
\subsection{The Ground State}
\label{sec:org4534998}

The ground state of the Medium is the condition of maximum undifferentiation —
fully random stochastic potential, no structure, no pattern, no resolved
distinction, maximum entropy. This is the overwhelmingly dominant condition of
the Medium. All of the structure we observe is a vanishingly small region of
differentiation — and vanishingly low entropy — within an infinite substrate that
is almost entirely at its stochastic ground state.

Order is the aberration. The stochastic ground is the norm.
\subsection{On Time}
\label{sec:orgabcae96}

Time does not exist in the Medium framework as a fundamental property. The Medium
has no past, present, or future. It simply is. What we experience as the passage
of time is the experience of transition — the change from one state of
differentiation to another, indexed by \(t\). The arrow of time is the arrow of
increasing entropy. Both are emergent properties of how differentiated structure
evolves through transitions. They are real as experiences. They are not
fundamental properties of the stochastic substrate beneath them.
\subsection{On Existing Physics}
\label{sec:org0b6fd7e}

Once locality has been established through structure formation, the existing
mathematical frameworks of physics are expected to emerge as valid descriptions
of differentiated structure behavior within their respective domains. They are
not wrong. They are downstream. The Medium framework does not replace them. It
describes the stochastic substrate from which the conditions that make them
applicable emerge.
\section{The Convergence With Human Thought}
\label{sec:orgd9a0190}

The framework arrived at through disciplined reduction converges with ideas that
have appeared independently across the deepest traditions of human philosophical
and scientific inquiry.

Across philosophical traditions and contemplative practice, independent inquiry
has consistently arrived at a structurally identical position: beneath the
apparent multiplicity of the world is one undifferentiated thing, and the
separateness of objects and beings is real as experience but not fundamental as
reality. The Medium framework arrives at the same place not through faith or
intuition but through the discipline of asking what the primitive must be and
rejecting everything that required something more primitive to explain it.

Within the physical sciences, independent lines of inquiry have converged on the
recognition that the ground of reality is probabilistic rather than deterministic
— that randomness is not a failure of description but a feature of the substrate.
The relationship between information, entropy, and the structure of physical
reality has been approached from multiple independent directions, each arriving
at the same formal identity between the counting of arrangements and the measure
of information.

The stochastic character of the Medium's ground state, the emergent nature of
time and locality, and the entropic arc from first distinction back to the ground
state — each of these has been approached from inside existing frameworks by
independent investigators working in different disciplines, arriving at
structurally consistent conclusions through entirely different methods.

That independent paths — physical, mathematical, philosophical, and contemplative
— arrive at the same destination is noted as significant. It is not offered as
proof. It is offered as a fact worth considering.

A subsequent version of this document will engage these convergences formally,
with full citations to primary and secondary sources, once those sources have
been studied directly by the author.
\section{Unresolved Questions}
\label{sec:org0c8e1eb}

These are the frontier of the framework. Each is stated precisely and carried
explicitly as an object of investigation.

\textbf{The Nature of Intrinsic Distinction.} What is \(\iota\)? The provisional
characterization offered in Section IV.V — as the property permitting local,
spontaneous entropy reduction in a fully random maximum entropy substrate —
connects \(\iota\) to thermodynamics and the stochastic framing without
circularity. Whether this characterization is formally adequate or requires
genuinely new mathematical formalism is the central open question.

\textbf{The Propagation Limit.} Is \(c\) a genuine primitive of the Medium or does it
emerge from \(\iota\)? If emergent, what appears to observers as a propagation
limit is a consequence of the structure of the observer rather than a constraint
of the substrate itself. The resolution of this question will significantly
affect the formal development.

\textbf{The Structure That Emerges.} What precisely is the self-sustaining structure
that forms when the structure formation condition is met? How does it form,
stabilize, and interact with other such structures? The formal definition of
these structures is the primary task of the next stage of development.

\textbf{The Stochastic Distribution.} The ground state is characterized as fully random
— no distribution shape, no weighted outcomes. How does the manifestation
probability \(p(x)\) relate to this characterization? If the ground state is
distribution-free, what governs \(p(x)\)? This tension between the formal
probability primitive and the fully random characterization requires resolution.

\textbf{The Observer.} The observer is herself a differentiated structure within the
Medium attempting to describe the Medium from within it. Every formal system
contains truths it cannot prove from within itself. The Medium framework faces
an analogous boundary, and acknowledges it.

\textbf{Consciousness.} Whether consciousness emerges from sufficient complexity of
differentiated structure is noted as significant and explicitly deferred. If it
belongs in the framework it must be derived — arrived at through formal
development, not assumed at the foundation.
\section{Summary of Axioms and Primitives}
\label{sec:orgb38b6d0}

\begin{center}
\begin{tabular}{lll}
Symbol & Definition & Status\\
\hline
Stochastic ground & Fully random autonomous, no distribution & New v0.2\\
\(\iota\) & Intrinsic Distinction — eternal, irreducible & Axiomatic\\
\(V\) & Activation of \(\iota\) — stochastic event & Consequence Ax. II\\
\(\mathcal{M}\) & Eternal Medium — no beginning, no end & Axiomatic\\
Axiom III & Locality Emergence — location produced at \(c\) & Axiomatic\\
\(M\) & The Medium — post-locality approximation & Provisional\\
\(\phi : M \rightarrow V\) & Population function — partial & Current\\
\(p(x)\) & Manifestation probability & Current\\
\(t\) & Transition index — not time & Current\\
\(M_t \rightarrow M_{t+1}\) & Transition between states & Current\\
\(c\) & Propagation limit & Provisional\\
\$ & \(\phi\)\textsubscript{t+1}(x) - \(\phi\)\textsubscript{t}(x) & \(\le\) c\$ & Propagation constraint & Provisional\\
\$ & \(\phi\)\textsubscript{t+1}(x) - \(\phi\)\textsubscript{t}(x) & = c\$ & Structure formation — origin of locality & Provisional\\
\(S = k \log W\) & Boltzmann entropy — ground state is \(S_{max}\) & New v0.2\\
\(\iota \Rightarrow S_{min}\) & First distinction is minimum local entropy & New v0.2\\
\(\iota \Rightarrow I_{max}\) & First distinction is maximum information & New v0.2\\
\end{tabular}
\end{center}
\section{A Personal Note}
\label{sec:org5f7002c}

I am seventy-eight years old. I am not an academic. I have no credentials in
physics or mathematics. I began this work with a question that would not leave
me alone and a conviction that the question deserved a serious attempt at an
answer.

I began this work knowing I may not live to see it reach its conclusions. That
does not trouble me. What matters is that the ideas are documented, attributed,
and made available to anyone who finds them worthy of development. If this
framework contains something true it belongs to whoever can take it further.
My name attached to it is enough — for the record, and for my family.

Gregor Mendel published his work on inheritance and it sat unread for decades
before being recognized as foundational. The ideas belong to whoever finds them.
The record belongs to their origin.
\section{License and Attribution}
\label{sec:org1ace46f}

This work is licensed under the Creative Commons Attribution 4.0 International
License (CC BY 4.0). You are free to share and adapt this material for any
purpose, including commercially, provided you give appropriate credit to the
original author, provide a link to the license, and indicate if changes were made.

\url{https://creativecommons.org/licenses/by/4.0/}

All versions of this document are archived on Zenodo and remain linked to the
authorship record of Bradley K. DePuy, independent researcher.

DOI: 10.5281/zenodo.18804587
\end{document}
